\section{Tail expansions and prescribed growth}
\label{ell:tail-calculus}

Bounded holomorphy alone does not allow the application of
Theorem~\ref{thm:path-selection}.  We prove the all-order expansion
and common-domain assertions needed there.  All expansions below
are finite to each requested precision.  An infinite formal
power-logarithm sum is never used as an inverse or a solution.

\begin{definition}[Tail expansion to finite precision]
\label{ell:tail-definition}
On an opened joined-type end, a field has an integer tail if, for
each nonintegral rate $B$, it can be written
\begin{equation}
 Y(\tau,\theta)=
 \sum_{l<B}e^{-l\tau}\sum_{j=0}^{d_l}\tau^jY_{lj}(\theta)
       +Y_{>B}.
 \label{ell:tail-series}
\end{equation}
Here $l\in\Z$.  There is a lower bound for $l$, each $Y_{lj}$ is a finite
equivariant spherical harmonic sum, and the sum is finite.
The remainder belongs to the cylindrical Sobolev space with any
fixed weight strictly below $B$.  A bound to finite precision
controls the finitely many displayed coefficients and that
remainder norm.  Polynomial contour costs can be paid by a
strictly smaller remainder rate.  For a growing field this is
a bound after subtraction of the finite nondecaying tail; it is
not a bound on the unsubtracted field at infinity.
\end{definition}

The extraction of homogeneous exponential terms when a weight crosses
an indicial root is part of the classical change-of-weight method
\cite[Section~5]{LockhartMcOwen1985}. Here the componentwise flat
operators allow an explicit polynomial calculation, including the
resonant logarithmic terms.

\begin{lemma}[Constant-operator tail improvement]
\label{ell:flat-tail}
Let $\mathcal P$ be either flat operator above.  Suppose a field
has a finite growth bound and its source has an expansion to
rate $B$, where the intermediate remainder weights avoid the
indicial roots.  Then the field has an expansion to that rate,
with free homogeneous coefficients recorded at the crossed
roots.  Estimates involve the source coefficients and remainder,
a fixed inner trace, and the previous growth bound.  They are
uniform on the contours of Lemma~\ref{ell:complex-inverse}, with
an arbitrarily small loss in a strictly available exponential
rate.
\end{lemma}

\begin{proof}
In one angular channel, apply
$(\partial_\tau-\kappa_1)(\partial_\tau-\kappa_2)$ to
$e^{-l\tau}p(\tau)$.  It gives $e^{-l\tau}$ times a constant
coefficient differential operator on the polynomial $p$.
If $-l$ is not a root, this is invertible on polynomials of a
given maximum degree by a triangular coefficient calculation.
If $-l$ is a root, integrate one polynomial factor, increasing
the degree by one; prescribe zero constant term for that
primitive and record the omitted constant as a free homogeneous
coefficient.  The roots are distinct, so this exhausts resonance.
Only finitely many angular channels occur in the forced finite sum.

Subtract these primitives.  Extend the remaining forcing by a
fixed cutoff to a full-cylinder source and solve it by the root
kernels at a new noncritical weight.  The difference from the
original field is homogeneous on the tail.  The previous growth
bound excludes roots growing faster than it permits; the roots
between the old and new weights are recorded explicitly.  The
remaining faster-decaying homogeneous modes are bounded by
the inner trace, using the exponential root gap and summing
their squared coefficients.  This proves both the improvement
and its estimate.  The kernel calculation on contours is the
one in Lemma~\ref{ell:complex-inverse}; polynomial losses are
absorbed by decreasing the remainder weight within its strict
gap.  No parameter-dependent global inverse at the new weight
has been invoked.
\end{proof}

\begin{lemma}[Nonlinear tails and finite prescribed growth]
\label{ell:prescribed-tails}
The small solution of Lemma~\ref{ell:analytic-family} on every
partial disconnection has integer tails with only $l>0$.
Its augmented linearization, denoted $\mathcal A_U$, can be
solved for a source with a finite-growth integer tail, with
arbitrary finitely supported prescriptions of free homogeneous
coefficients in the nondecaying indicial channels. Unspecified
coefficients are zero, so the prescribed growth has a finite maximum.
The result has an integer tail; its growth does not exceed the
maximum of zero, the forced growth, and the prescribed growth,
apart from polynomial factors.  The solution is unique with
these prescriptions and the other augmented data.  Estimates
to any fixed finite precision are linear in the data, with
constants allowed to depend on that precision and growth.
Every finite tail coefficient depends analytically on the
parameters on one common domain.
\end{lemma}

\begin{proof}
First consider the small nonlinear solution.  Its initial
positive-weight estimate and
Equation~\eqref{ell:flat-equation} give a source at twice a
slightly weaker positive decay rate.  Apply
Lemma~\ref{ell:flat-tail}.  Suppose finitely many powers have
already been extracted and the remainder has positive rate $B$.
In the difference between the nonlinear expression and its
finite Taylor expansion, every occurrence of that remainder
is multiplied by a field or derivative with a fixed positive
decay rate.  The analytic Taylor remainder is controlled by
Sobolev multiplication, since $H^k$ is an algebra at the chosen
orders.  Consequently the next source remainder has rate at
least $B+\epsilon$ for one fixed sufficiently small
$\epsilon>0$, using intervening noncritical weights when needed.
This improves indefinitely.  The initial roots are integers;
Taylor products add their degrees.  The angular property in
Lemma~\ref{ell:analytic-family} makes each resulting coefficient
a finite harmonic sum.  The original positive decay excludes
all nonpositive degrees.  This proves the first assertion,
including estimates on the complex contours.

On an opened end the linearization is
$\mathcal A_U=\mathcal P+\mathcal B_U$, where the coefficients
of the differential operator $\mathcal B_U$ have positive
integer tails and decay.  Start at the most growing required
power of the desired solution.  Use the polynomial primitives
in Lemma~\ref{ell:flat-tail} to cancel it.  Multiplication by
$\mathcal B_U$ improves the degree by a positive integer, so
the recursion is triangular in degree.  At each indicial
channel fix the log-free homogeneous coefficient to be its
prescribed value; use the zero-constant polynomial primitive
for the forced part.  Continue through finitely many degrees
until the residual decays strictly faster than the basic
weight $\nu$.  This constructs a finite sum $T$, with no
growth greater than the allowed maximum.

Choose a fixed cutoff $\zeta$ supported on the opened end and
equal to one sufficiently far out.  With all finite seam data,
core tests, and compact source unknowns included in
$\mathcal A_U$, set
\begin{equation}
 Y=\zeta T+
       \mathcal A_U^{-1}\bigl(f-\mathcal A_U(\zeta T)\bigr).
 \label{ell:finite-growth-inverse}
\end{equation}
The term in parentheses is in the basic decaying source space.
The inverse is exactly the inverse already constructed in
Lemma~\ref{ell:analytic-family}.  No smallness of $T$ or of this
source is required: this is a linear solve.  The correction
decays and cannot change any prescribed nondecaying coefficient.
Applying Lemma~\ref{ell:flat-tail} to its equation, and using
the coefficient tails of $\mathcal B_U$, obtains all further
orders.  If two solutions have the same data, their difference
has no nondecaying tail, hence belongs to the basic decaying
space and vanishes by its inverse.  This proves uniqueness and
the growth assertion.

All steps in Equation~\eqref{ell:finite-growth-inverse} are
finite polynomial operations, fixed cutoffs, and the same
analytic inverse family.  A leading coefficient can alternatively
be extracted by multiplying the tail by its leading exponential
and removing its known polynomial terms, then taking its limit
at infinity.  The remainder estimates are locally uniform in
the parameters, so these limits are holomorphic.  Repeat only
finitely many times for any requested coefficient.  The
parameter domain never shrinks as the precision or prescribed
growth increases; only estimates and a convenient tail-start
threshold may change.  In particular, a cutoff moved farther
out on an opened infinite end does not change the admissible
range of another edge's finite length.
\end{proof}

\begin{lemma}[Triangular auxiliary equations]
\label{ell:triangular}
Lemma~\ref{ell:prescribed-tails} applies successively to every
finite system appended to the small solution in which each
equation has diagonal operator $\mathcal A_U$ and source given
by Taylor differential operations on earlier fields.  Allow
finite growth at already opened ends, and free asymptotic
conditions depending polynomially on finitely many earlier
tail coefficients.  On still-finite seams impose homogeneous
matching and zero added critical tests.  Require that sources
on those halves have at least two decaying factors, counting
the base field and its derivatives.  On opening a new pair
of ends, start with decay there.  The resulting fields and
the inverse of the triangular linearization have uniform
finite-precision bounds and the same analytic parameter domain.
\end{lemma}

\begin{proof}
Solve the equations in their given order by
Equation~\eqref{ell:finite-growth-inverse}.  Products of finitely
many fields with finite growth still have finite growth.
Forced nondecaying coefficients are determined by the equation;
only the free indicial constants are prescribed.  Thus a forced
power is never incorrectly set to zero by an asymptotic boundary
condition.  Conditions involving earlier coefficients are known
data at the relevant step.  On a finite bent end the two
decaying factors supply the exponential surplus needed to pay
the polynomial strengthening on that same end.  Linearizing
the finite triangular system leaves $\mathcal A_U$ on every
diagonal, so induction gives its inverse as well.

At a newly opened end every earlier field initially decays.
The off-diagonal coefficients of this linearization therefore
decay there.  This fact will preserve the leading incoming
mode calculation when this end is closed again.  On older
opened ends the off-diagonal coefficients can grow, but their
growth is finite and is handled at each successive step.
Arbitrarily large bounds for a finite auxiliary field do not
affect the smallness of the base metric or the domain of its
inverse.  This proves all the assertions.
\end{proof}

